\section{Research Gaps and Directions}
\label{sec:discussion}
\paragraph{Measuring non-target preservation.}
Image-quality metrics are insensitive to many design-specific failures: a corrected headline can shift a logo, modify a brand color, or alter the line breaks of an unrelated paragraph. Editing benchmarks should include protected-object masks, semantic content checks, and explicit layout tolerances. Preservation must be evaluated on the correct level: pixels for raster methods, paths for SVG, object identities for layered documents, and DOM/content structure for code-native artifacts. However, some non-target movement is necessary: InfoEdit demonstrates that infographic editing may require coordinated global reflow \citep{yang2026infoedit}. An evaluation protocol should distinguish permitted dependent changes from changes to explicitly protected content (Eq.~\ref{eq:reflow}).

\paragraph{Native artifact validity.}
An image preview cannot certify that a PDF, SVG, PPTX, PSD, or HTML/CSS design remains editable. Reliable evaluation should include a round-trip: reopen the saved result in a target editor or renderer, inspect whether objects retain the expected type, perform a controlled edit, then render again. This is especially important for code-native agents that may output visually convincing pages with brittle styles or missing assets \citep{designascode2026,editablevisual2026}.

\paragraph{Representation-aware comparisons.}
A fair comparison must define the task and permitted output. Penalizing raster editing for not generating vector nodes, or rewarding any SVG solely because it contains paths, mixes two evaluation questions. We recommend separate leaderboards for semantic edit compliance, format validity, and visual result quality, with clear reporting of which dimensions are inapplicable to a given output contract.

\paragraph{Human-centred outcomes.}
Designer workflows involve more than one-shot aesthetic ratings. Vinci provides an example of measuring editing efficiency \citep{vinci2021}. Future comparisons should additionally test task completion time, correction burden, how easily users can localize an error, compatibility with editing software, and designer trust across multiple changes. User expertise and task familiarity should be treated as experimental variables.

\paragraph{Reliable evidence synthesis.}
Evaluation descriptions found in abstracts should not automatically be mapped to specific numerical scores. Full experimental tables and appendix details can reveal sample-size limits, unreported exclusions, or test conditions that change the interpretation. The remaining survey work therefore includes systematically logged searches and source audits, explicit boundary adjudication, bibliographic reconciliation, and separate flags for measured versus merely claimed editing capability.

\paragraph{Reproducibility and provenance.}
Source code, weights, evaluation data, and proprietary API dependencies introduce independent reproducibility constraints. A system may publicly release an inference pipeline without the training recipe or may require a paid closed model at evaluation time. These attributes should be tracked separately from whether a final artifact is editable, and should not be inferred merely from an available project-page link.
